\chapter{Referencial teórico}


\section{Negócios agroflorestais sustentáveis}

A crescente preocupação com a crise ambiental global, como resultado do uso excessivo dos recursos naturais da Terra e do aumento populacional \cite{Steffen2015PlanetaryPlanet}, tem exigido mudanças fundamentais nos sistemas de consumo e produção. O desmatamento é um grande contribuinte para as emissões globais de gases de efeito estufa, representando 18\% dessas emissões \cite{ONU2015}. O desmatamento causa diversos efeitos ambientais negativos, tais como perda de biodiversidade \cite{Pyles2018LossHistories}, deslocamento de encostas e erosão \cite{Holanda2021SoilOverview}, insegurança alimentar  \cite{TataNgome2019AssessingCameroon}, e desalojamento forçado de comunidades tradicionais  \cite{Ingalls2018TheRegion}. A pressão ambiental também afeta a Floresta Amazônica, que sofre forte pressão com a conversão das florestas em pastagens e terras agrícolas \cite{SEEG2021SystemSEEG}.
As empresas desempenham um papel fundamental na preservação dos recursos naturais e da biodiversidade das florestas tropicais. Atualmente, muitas empresas estão criando produtos verdes ou adicionando valor às suas marcas por meio de iniciativas ecológicas \cite{Majid2015GivingMarket}. No entanto, muitas vezes, essas iniciativas são limitadas a pequenas empresas que enfrentam dificuldades para agregar valor aos seus produtos e comercializá-los sem a intermediação de atravessadores \cite{Caniato2012EnvironmentalResearch}.

A relevância da produção agroflorestal para o mundo é indiscutível, por  ela que alimentos e matérias-primas são fornecidas a todos os setores da economia. No âmbito internacional da FAO afirma-se que a agricultura familiar é fundamental para garantir segurança alimentar e proteção ao meio ambiente, entre outros benefícios \cite{FAO2020}. Esse grupo de empreendedores, dada as suas peculiaridades, necessitam de apoio de políticas públicas para se tornarem ainda mais produtivos e sustentáveis. 

Definir o desenvolvimento agroflorestal sustentável requer um esforço observacional e prático ao desenvolvimento rural sustentável, pois, este ambiente vem sofrendo profundas transformações em suas demandas e necessidades. O desenvolvimento que antes se apresentava majoritariamente como produção de subsistência, hoje dá lugar a um complexo sistema agroindustrial \cite{Bastos2018DeterminantesBrasileiro} e social. É importante neste sentindo compreender que definir o desenvolvimento rural com apenas um conceito seria uma proposição simplista do contexto de crescimento no meio rural. Partindo da definição de consequência de ações governamentais definidas por \citeonline{Navarro2001DesenvolvimentoFuturo} como “atitudes práticas”, este autor descreve que:

\begin{citacao}
Desenvolvimento rural e florestal, portanto, pode ser analisado a posteriori, neste caso se referindo às análises sobre programas já realizados pelo Estado (em seus diferentes níveis) visando a alterar facetas do mundo rural a partir de objetivos previamente definidos. Porém pode se referir também à elaboração de uma ação prática.
\end{citacao}

Neste campo econômico, o sistema agroflorestal sustentável, especialmente a agricultura de base familiar e silvicultura de base comunitária, são componentes indispensáveis do ponto de vista econômico, social e ambiental. Mesmo no mundo de hoje, que enfrenta desafios como superpopulação, mudanças climáticas, dificuldades na produção agrícola em grande escala e competição de recursos, esta atividade desempenha um papel único \cite{Riegler2018,Rockstrom2017}.

Ao atendar aos requisitos do desenvolvimento sustentável, a agricultura familiar busca eliminar a fome no meio rural e ajuda os agricultores a gerarem renda, oportunidades de trabalho e a possibilidade de competir com os grandes produtores a partir do cooperativismo e o associativismo, a fim de fornecer alimentos mais saudáveis por culturas mais diversificadas, em que os produtores rurais cultivam várias safras na mesma terra \cite{Santana2014}.

Segundo a \textit{Food and Agriculture Organization of the United Nations} FAO \cite{FAO2017a} a agricultura para consumo inclui todas as atividades agrícolas e familiares e está ligada a várias áreas do desenvolvimento rural, seja a produção de hortaliças e alimentos ou de carne para consumo alimentar. A agricultura familiar é uma forma de organizar a produção agrícola, florestal, pesqueira, pastoril e aquícola gerida essencialmente por familiares.

Ademais, no que diz respeito à silvicultura sustentável, principalmente à de base comunitária, incluem-se iniciativas, ciências, políticas, instituições e processos que visam aumentar o papel da população local no governo e gestão dos recursos florestais. Ela abrange tanto os regimes de manejo florestal comunitário quanto a silvicultura de pequenos proprietários \cite{Sabogal2016Community-basedForestry}. O Manejo florestal comunitário é um conceito de silvicultura de base comunitária desenvolvido por comunidades locais, sejam elas tradicionais ou ali assentadas \cite{Massiri2020InstitutionalIndonesia}. Este tipo de manejo envolve a comunidade na governança florestal, podendo se organizar de forma autônoma e desenvolver instituições locais para regular o uso dos recursos naturais e gerenciá-los de forma sustentável, tanto para o extrativismo sustentável de seus recursos quanto para a manutenção e recuperação da florestal em que são alojados \cite{Gilmour2016FortyEffectiveness}.

As Unidades de Manejo Florestal comunitárias (UMFC), desempenha papéis significativos nas práticas de manejo florestal sustentável em todo o mundo \cite{Osie2020Trade-offsManagement}. As práticas de manejo florestal realizadas pela UMFC não são apenas para garantir a sustentabilidade do ecossistema, mas também para acomodar os interesses da comunidade local. Comunidades locais que vivem em torno de áreas florestais têm um grande interesse nos recursos florestais e o manejo correto destes recursos salvaguarda a floresta em pé.

Além do manejo florestal desenvolvido pelas unidades de manejo florestal, outro fator de suma importância para a efetiva sustentabilidade florestal é a compreensão  da agricultura familiar florestal no Brasil. A agricultura familiar florestal é outro modelo econômico de base social para populações envolvidas no cenário floresta. \citeonline{Mudge2014FarmingLivros} a define como um sistema de cultivo intencional de produtos florestais não madeireiros (PFNMs) sob um dossel florestal, tendo como base a agricultura, silvicultura, pesca, sistema silvipastoril e produção aquícola, promovida como uma forma de conservar espécies florestais ameaçadas. Estas unidades agrícolas são geridas e operadas por uma família predominantemente dependente do trabalho familiar, incluindo mulheres e homens. Nestes sistemas, a família e a fazenda estão ligadas e coevoluem em: economia, meio ambiente, funções sociais e culturais \cite{FAO2014}.

\section{Empreendedorismo para negócios sustentáveis e novos paradigmas}

O empreendedorismo, quando relacionado à criação de negócios, não consiste apenas na criação de uma empresa, ser empreendedor vai além de ter seu próprio negócio e implica em uma visão de mundo diferente, uma mudança de paradigma, de pensamento. \citeonline{Garcia2012} define o empreendedorismo inovativo como sendo a habilidade de criar e de construir algo a partir do nada, sendo, desta forma, um ato humano nascido da criatividade. O empreendedor é um constante inovador, sendo aquele que está sempre em busca de soluções, que consegue enxergar nas oportunidades, tendo iniciativa e sendo proativo e visionário \cite{Milian2020}.

Dessa forma, ele pode contribuir para introduzir inovações na empresa em que trabalha, além de ajudar a solucionar problemas da sua cidade, entre outras possibilidades que se fazem presentes desde um comportamento empreendedor é desenvolvido \cite{Alencar2019,Loiola2016}. O alcance do sucesso ao empreender um novo negócio  em todas as suas vertentes é um fenômeno complexo e objeto de discussões e pesquisas no âmbito acadêmico e profissional \cite{DEVECE20165366,Calic2016}. De acordo com \citeonline{DeMori1998}, define como fatores condicionantes do sucesso no empreendedorismo: características individuais do fundador, características estruturais e estratégicas do novo negócio, e condições determinadas pelo meio ambiente da empresa.

Sob essa perspectiva, o empreendedorismo com foco sustentável é uma forma de impactar o desenvolvimento sustentável para além do âmbito econômico, mas com impactos social e ambiental, por meio do processo de descobrir novas oportunidades de negócios inovadores e sustentáveis com os seus enfoques distintos \cite{Sepulveda2011,Davidsson2016}. O empreendedorismo com foco na sustentabilidade pode maximizar lucros de oportunidades de negócios, contribuindo para redução da pobreza e impactando positivamente na sociedade e no meio ambiente \cite{Medauar2020}.

Considerando aspectos relacionados à proteção ambiental, preocupação com os recursos naturais e consumo, e redução da pobreza — o objetivo dos empreendedores é incorporar práticas sustentáveis na criação de negócios — buscando a promoção e o desenvolvimento sustentável em prol da sociedade e o meio ambiente \cite{raufflet2014}

Atender às necessidades ambientais sem a contrapartida da onerosidade é uma tarefa muitas vezes difícil de ser realizada, principalmente para as pequenas empresas. Buscando contrabalancear este entrave, atualmente as empresas estão reforçando suas estratégias de marketing, já que a aplicação de técnicas de criação de mídia vem se mostrando um importante fator para o sucesso do posicionamento estratégico de empresas, independentemente da natureza ou escala de sua atividade. Uma forma de demonstrar que a empresa sustentável explora as necessidades dos consumidores é através do gerenciamento estratégico de sua marca \cite{Kotler2019MarketingManagement}.

\textit{Branding} é um processo de criação de nomes ou símbolos para diferenciar um produto de seus concorrentes \cite{Kotler2019MarketingManagement}. Quando bem utilizado, o \textit{Branding} é o ator proativo de dar vida à empresa, utilizando esforços concentrados e coordenados de todas as principais partes interessadas na operacionalização e gerenciar sua marca \cite{Marland2017GovernanceBranding}. Quando empregado em pequenas empresas, o \textit{Branding} é o controle estratégico de elementos que comunicam sobre todas as capacidades competitivas da empresa, mostrando ao consumidor que seu potencial pode ser parelho ao dos seus concorrentes de maior porte. As técnicas de \textit{Branding} incluem publicidade, embalagem e design de produto \cite{Ho2017DetanglingSuccess}. 

O marketing com foco no \textit{Branding} pode ser usado para criar demanda para um produto ou serviço, a exemplo as técnicas de \textit{Product Naming}\cite{Bergh2016SoundNames}. Isso pode ser alcançado criando uma aparência única para a marca e destacando seus recursos e benefícios por meio de mídias. Além disso, concentrando-se nas necessidades do cliente, as empresas podem criar campanhas de marketing personalizadas para seu público-alvo, utilizando técnicas de psicologia do consumidor \cite{Schmitt_Be2012}. Ao entender os desejos e necessidades dos clientes, as empresas podem criar estratégias para garantir o desenvolvimento de seus produtos, considerando potenciais acréscimos em sua produção em consequência das práticas sustentáveis \cite{Rubio2021}.

O feedback pode ser usado para ajustar as campanhas com base no comportamento e nas emoções do cliente \cite{Zavattaro2021PublicDevelopments}. Essa abordagem ajudará a  Inc. a criar campanhas de marketing mais personalizadas para o cliente, o que pode aumentar o engajamento e o \textit{Brand Loyalty}\footnote{O \textit{Brand Loyalty} descreve os sentimentos positivos dos consumidores sobre uma marca e sua dedicação em comprar repetidamente os produtos ou serviços dessa marca, independentemente de deficiências, ações dos concorrentes ou mudanças no ambiente \cite{Keller2006BrandsPriorities}.}, fatores-chave para a construção do \textit{Brand Equity} \cite{Keller-Aviram202187}.



\section{Marcas como estratégia de marketing}


Contudo, o estabelecimento de empreendimentos florestais sustentáveis necessita de uma abordagem gerencial distinta, visto que os processos também são consideravelmente distintos e frequentemente mais dispendiosos do que os métodos convencionais \cite{Al-Minhas2020431}. Para alcançar um desenvolvimento sustentável e uma posição competitiva elevada, as empresas das regiões estão procurando a integração de suas políticas ambientais às suas estratégias de marketing, através do marketing ambiental \cite{Conejero20126}. O campo emergente do empreendedorismo socioambiental busca preencher essa lacuna e promover o desenvolvimento sustentável. Um aspecto chave desta estratégia é a gestão sustentável das florestas, que requer uma abordagem de gestão distinta e um investimento significativo.

O marketing e as estratégias de negócios baseados na Sustentabilidade e Tecnologias Apropriadas são importantes na determinação de quais tecnologias devem ser priorizadas para investimento, tanto para resolver questões ambientais básicas do país quanto para ampliar a exportação de produtos para países que valorizam a preservação do meio ambiente \cite{Medrano2020TheOrientation}. Dentre as ferramentas de amplo uso pelo marketing está a marca.

No campo do marketing sustentável, a marca é considerada um ativo, ao representar a identidade e a reputação de uma empresa perante seu público. É através da construção de uma marca forte e reconhecida que uma empresa alcança a diferenciação, de seus concorrentes e conquistar a lealdade de seus clientes, assim como cria a fidelização dos consumidores perante seus produtos. Além disso, a construção de uma marca forte é fundamental para o sucesso a longo prazo da empresa.

A marca é composta por diversos elementos, como logotipo, slogan, cores, nome entre outros, que juntos formam sua personalidade, seu \textit{brand equity}. É importante que esses elementos sejam consistentes e coerentes com a mensagem que a empresa deseja transmitir, isso é que a marca seja a porta-voz da empresa perante seus consumidores. Além disso, a marca também é construída através das experiências e percepções dos clientes com a empresa, seja por meio de sua comunicação, produtos e serviços, atendimento ao cliente, entre outros aspectos.

A grande importância empresarial das marcas necessita de um sistema gerencial específico, sistema esse conhecido como branding. Investir em uma boa estratégia de branding é fundamental para o sucesso de empresas de manejo florestal sustentável, ao permitir que ela se destaque no mercado e conquiste a confiança e lealdade de seus clientes \cite{Aprile2022HowItaly}. Uma marca forte e bem estabelecida pode ser um diferencial competitivo importante ao permitir que a empresa tenha uma vantagem sobre seus concorrentes e consiga cobrar preços mais elevados por seus produtos ou serviços. 

Atender às necessidades dos consumidores através da aplicação de técnicas \textit{Branding} é um dos fatores-chave para o sucesso do posicionamento estratégico de uma marca, independentemente da natureza ou escala da sua atividade \cite{Chang2018EnhancingMarketing}. As novas demandas de informação imbuíram às empresas a necessidade de um gerenciamento de marketing bem definido, haja vista os novos padrões de mídia e a possibilidade de entrada de novos competidores parelhos aos demais do seu mercado. Neste sentido, as empresas agrícolas não ficam de fora. Embora as empresas tenham muitas peculiaridades e diferenças em termos de escala e flexibilidade, a aplicação do \textit{Branding} para marcas é um conceito universal e, desta forma, passível de adaptação e implementação em qualquer natureza negocial. 

A marca é um signo distintivo de um produto, empresa ou serviço, visto que possuem amparo legal albergado na formalização de pedido de registro junto a organismos regulamentadores, geralmente designados por entes públicos de cada nação \cite{Sammut-Bonnici2015BrandBrandingb}. Em função disso, torna-se ativo econômico para os seus requerentes. Constitui uma subárea do direito de Propriedade Industrial, definida como “conjunto de referenciais físicos e simbólicos capazes de influenciar e determinar a preferência para os produtos, havendo como base a oferta de valor a ela associada” \cite{Mourad2014}. Tendo validade em âmbito nacional, um registro de marca é válido por um período de dez anos. Contudo, a pedido do titular, pode ser prorrogado várias vezes, por iguais e sucessivos períodos. Conforme a legislação vigente no Manual de Marcas \cite{INPI2017} esses sinais distintivos são conceituados sob três aspectos, sendo: (i) natureza (marca de produto, serviço, coletiva e certificação), (ii) formas de apresentação (normativa ou verbal, figurativa, mista e tridimensional), e (iii) princípios legais (territorialidade, especialidade e sistema atributivo).

A Propriedade Intelectual possui um conceito amplo que introduz temáticas relacionadas ao processo de proteção às criações da mente, envolvendo aspectos relacionados à Propriedade Industrial, Direito Autoral e Direitos Conexos. Observa-se que há uma preocupação em relação à proteção da referida propriedade, como constatado nos principais tratados e acordos a seguir: Convenção da União de Paris (CUP), Tratado de Madri, Convenção de Berna, Convenção de Novas Variedades de Plantas, Acordo Geral de Tarifa e Comércio, Acordo TRIPS (\textit{Trade Related Aspects of Intellectual Property Rights}) e a Rodada Uruguai \cite{FERRAZ2008}.

O Instituto Nacional da Propriedade Industrial (INPI), criado em 1970, é uma autarquia federal vinculada ao Ministério da Indústria, Comércio Exterior e Serviços (MDIC), responsável pela execução das normas que regulamentam a propriedade intelectual para indústria no Brasil em âmbito social, econômico, jurídico e técnico, cabendo ainda o pronunciamento quanto ao interesse de assinatura, ratificação e denúncia de convenções, tratado, convênios, assim como acordos sobre a referida propriedade \cite{INPI2020}, destacando-se a “marca” como relevante fator para distinguir um produto ou serviço entre organizações. No plano internacional, a Organização Mundial de Propriedade Intelectual (WIPO) é uma das agências especializadas na proteção global da propriedade intelectual. Com sede em Genebra, é composta atualmente de 187 estados-membros e controla, sob sua tutela, 27 tratados internacionais \cite{INPI2020}.

A marca, neste contexto, entra como ativo intangível capaz de formular os valores na mente dos consumidores, mantendo a empresa em posição competitiva parelha. Nesta linha, as pequenas empresas devem orientar-se a explorar os fatores ambientais imediatos e mais amplos passiveis de apresentação, transformando os compromissos social e ambiental em um indicador de responsabilidade ética empresarial, utilizando-se dos novos padrões de consumo e os novos estilos de vida \cite{Thelwall2021LifestylePatterns}.


Geralmente as marcas que visam preservar os recursos naturais e a biodiversidade de nosso planeta desenvolvem apenas a comercialização da produção de pequenas empresas oriundas de associações ou estrutura familiar. Este tipo de atividade comercial, quando dependente de atravessadores, geram incertezas e baixo valor agregado aos produtos \cite{Caniato2012EnvironmentalResearch}, já que as pequenas empresas estão presentes apenas na cadeia de suprimentos, impedidas de agregação de valor e posterior comercialização. O problema se agrava enquanto as grandes companhias se apegam apenas aos pressupostos de marcas verdes, esquecendo-se dos fatores sociais necessários e inerentes à produção sustentável. 

As estratégias de green branding têm sido amplamente adotadas pelas empresas como uma forma eficaz de criar uma imagem de marca positiva, aumentar a vantagem competitiva e melhorar o desempenho financeiro \cite{Chang2012,Visser2018664,Yusof2020GoingPerformance}. Estudos mostram que os consumidores preferem marcas ecologicamente corretas em vez de marcas que não agregam iniciativas sustentáveis \cite{Pimonenko2020}. A exemplo, consumidores costumam pagar mais por carros ecologicamente corretos \cite{GahlotSarkar2019BrandAppeals,Leckie2021PromotingBrands}, demonstrando uma forte intenção de sustentabilidade percebida pelo consumidor \cite{Lin2012DoubleUsage}.

Neste cenário de mudança de paradigma, o papel da sustentabilidade ambiental e os novos padrões de consumo surgem como influenciadores significativos no planejamento estratégico das marcas empresariais \cite{Chen2019TrueBehavior}. A percepção crescente do público em relação à importância da proteção ambiental está moldando um novo consumidor mais consciente e exigente, resultando na redefinição das práticas de compra e dos valores da marca \cite{Hatch2012TheExpression}. 

Assim, a marca forte, além de representar a identidade corporativa, deve incorporar uma visão de sustentabilidade e atuar como um catalisador na promoção de práticas de consumo responsáveis. Isso indica um redirecionamento para as empresas que buscam estabelecer uma marca resiliente, onde a sustentabilidade não é apenas um aditivo, mas um elemento essencial integrado ao núcleo da estratégia de branding. Nesse contexto, o investimento em inovação e práticas de negócios sustentáveis, aliado a uma marca sólida, pode impulsionar uma nova era de consumo sustentável e posicionar a empresa como líder em seu segmento, alinhada aos princípios do desenvolvimento sustentável \cite{Kotler2019MarketingManagement}.

No entanto, nem todas as empresas conseguem comunicar com sucesso suas práticas ou iniciativas verdes, aumentando assim a disparidade entre as grandes corporações que utilizam marcas vedes e as pequenas empresas que poderiam utilizar e difundir seus produtos com valor premium. Um motivo importante, mas frequentemente esquecido, é a assimetria de informações, tornando difícil para os consumidores entenderem a diferença entre uma marca verde de uma pequena empresa e quais são os benefícios associados \cite{Chi2021UnderstandingDestinations}. 

Embora a satisfação, confiança e lealdade do consumidor em relação às marcas tenham sido amplamente estudadas, essas questões ainda não foram discutidas sob a perspectiva do branding para pequenas empresas que desenvolvem atividades no contexto florestal \cite{Molinillo2019ASuccess}.  

Portanto, estudos teóricos e metodológicos que integrem antecedentes e resultados práticos da relação consumidor-marca verde são de extrema importância para o meio científico e social. Estudos nesta perspetiva visam enriquecer esse campo ao discutir o branding para pequenas empresas que atuam no contexto florestal sustentável.


\section{Sustentabilidade ambiental e novos padrões de consumo}

O desenvolvimento pleno da sustentabilidade ambiental implica necessariamente na proteção e melhoria das florestas atualmente existentes, proporcionando ações de restauração, manutenção e regulamentação ambiental a serem aplicadas às diferentes unidades de manejo. Antes, ações de sustentabilidade ambiental eram direcionadas somente para a elite econômica, negligenciando questões sociais. Agora, são vistas como cruciais para o desenvolvimento sustentável do país. 
 
As iniciativas para promoção de mudanças nos padrões de consumo da população podem ser caracterizadas em dois grandes grupos. O primeiro atrelado à política de modificação da demanda de alimentos. Esta frente de trabalho é ligada a políticas para modificar a oferta de alimentos saudáveis, em que se discutem os sistemas alimentares para estabelecer um ambiente adequado quanto aos impactos da alimentação sobre a saúde, com ingestão adequada de nutrientes, portanto ingestão de energia suficiente, sem excessos, voltada para atender às necessidades energéticas do indivíduo \cite{Maia2018}. Além de que tais alimentos sejam oriundos de ambientes sustentáveis, impactem minimamente a natureza e gerem menos sofrimento animal em sua produção. Nesta segunda vertente está a agricultura sustentável e a silvicultura sustentável. 
 
Segundo a \textit{Food and Agriculture Organization of the United Nations} FAO \cite{FAO2017a}, a agricultura para consumo inclui todas as atividades agrícolas e familiares e está ligada a várias áreas do desenvolvimento rural, seja a produção de hortaliças e alimentos, seja a produção de carne para consumo alimentar. A agricultura familiar é uma forma de organizar a produção agrícola, florestal, pesqueira, pastoril e aquícola gerida essencialmente por familiares.

É sabido que a rápida expansão da pandemia COVID-19, que envolveu 186 países entre dezembro de 2019 e março de 2020, agravou os riscos de insegurança alimentar entre grave e extrema de 135 milhões em janeiro de 2020 para 265 milhões até o final de 2020 \cite{Dongyu2020}. O grave problema da insegurança alimentar afetará as populações dos países desenvolvidos e em desenvolvimento \cite{Lal2020}. No âmbito internacional da FAO, afirma-se que a agricultura familiar é fundamental para garantir segurança alimentar e proteção ao meio ambiente, entre outros benefícios \cite{FAO2020}. Esse grupo de empreendedores, dadas as suas peculiaridades, necessitam de apoio de políticas públicas para se tornarem ainda mais produtivos e sustentáveis. Apesar de pelo menos 90\% das propriedades rurais do mundo serem um negócio familiar, não há uma definição universal para agricultura familiar.

Os diversos conceitos relacionam a propriedade e seu gerenciamento, o uso de mão-de-obra e tamanho físico ou econômico da citada. Existe um consenso de que a propriedade rural é operada ou gerida por um membro de um agregado familiar que o trabalho é realizado pelo proprietário e sua família \cite{Schneider2003}. 

Em 2018, o Instituto Brasileiro de Geografia e Estatística (IBGE), em parceria com o Ministério do Desenvolvimento Agrário (MDA), publicou a parcial do Censo Agropecuário 2017 \cite{IBGE2018CensoParanaense}, relatando que, no Brasil, atualmente 23\% das propriedades rurais são de agricultores familiares, ratificando a importância da Agricultura Familiar no Brasil. A agricultura familiar é fundamental para a provisão de alimentos-base para a constituição de dietas saudáveis, tais como legumes, frutas, vegetais e de origem animal. Apesar de os 4.367.902 de estabelecimentos familiares representarem 84,4\% do total, estas propriedades ocupam apenas 24,3\% da área ocupada pelos empreendimentos agropecuários brasileiros. 

Confirmando os estudos da FAO, a agricultura familiar é fundamental para a  segurança alimentar do Brasil, responsável por 87,0\% da produção nacional de mandioca; 70,0\% da produção de feijão, 21,0\% da produção do trigo e 16,0\% da soja, principais produtos da pauta de exportação brasileira. Essa modalidade de propriedade também é responsável por 59,0\% do plantel de suínos, 50,0\% do plantel de aves, 30,0\% dos bovinos, \cite{IBGE2018CensoParanaense}.

Nesta perspectiva, a literatura especializada reporta que o desenvolvimento econômico sustentável depende de uma integração qualificada de empresas, sustentabilidade ambiental e comunidades locais envolvidas \cite{Hu2021ResearchData}, considerando os povos originários e a população local agentes propagadores da sustentabilidade. Algumas das formas de fomentar estas iniciativas se dão por meio da silvicultura de base comunitária e a agricultura familiar \cite{Alvard1994ConservationPeoples,Mateo-Vega2017FullDarien}.



\section{Políticas públicas de apoio a Desenvolvimento de tecnologias ambientais sustentáveis}

Inicialmente a ação ambiental era dirigida à elite econômica, de modo que os ministérios responsáveis se distanciaram das demandas que permeavam o tecido social do país. A ação ambiental é hoje vista como um fator importante na promoção do desenvolvimento sustentado e sustentável de um país. Se, por um lado, a Revolução Industrial trouxe novos cenários para a economia e as atividades humanas, por outro, novos processos, tecnologias e produtos mudaram a percepção das pessoas sobre os recursos naturais, os seres humanos e o planeta. Alinhados com tais políticas ambientais estão os Objetivos de Desenvolvimento Sustentável (ODS) \cite{FAO2014}.

A agricultura e silvicultura sustentável de pequeno porte, sendo o lastro do consumo alimentar mundial, tem se mostrado uma saída para o enfrentamento da fome em todo o mundo. Segundo a FAO e a Organização Pan-Americana de Saúde/Organização Mundial de Saúde (OPAS/OMS) \cite{FAO2017}, organizações responsáveis por estudos no âmbito da segurança alimentar mundial, em relatório intitulado “Panorama da Segurança Alimentar e Nutricional na América Latina e no Caribe”, consoante ao Objetivo do Desenvolvimento Sustentável (ODS), visam a: acabar com a fome, alcançar a segurança alimentar, melhorar a nutrição e promover a agricultura sustentável.

Os ODS exigem o desenvolvimento de soluções multidisciplinares nas quais todos os campos de estudo relevantes se unem para criar um plano viável. Da agricultura à educação, todos os aspectos da sociedade precisam ser considerados ao lidar com a desnutrição e a fome. Isso envolve a criação de sistemas agrícolas sustentáveis que considerem as necessidades nutricionais. Também precisa ser considerado esses sistemas afetam o meio ambiente e outros aspectos da sociedade \cite{Silva2018g}.

Para o desenvolvimento sustentável, a ciência e a tecnologia correspondem a um sistema de articulação entre uma racionalidade ambiental do processo de desenvolvimento e os processos concretos que definem as possibilidades de estratégias de manejo integrado do meio ambiente. Essa interação requer que o sistema de ciência e tecnologia esteja sustentado por paradigmas que incorporem o potencial ecológico, as condições ambientais e os valores culturais na organização dos processos produtivos \cite{Furlan2018}.

A busca por inovação de processos e de produtos possibilita o alcance dos níveis mais elevados de competitividade. Para as empresas serem bem-sucedidas em mercados internacionais, são necessárias vantagens competitivas baseadas em inovação, complexidade e/ou sofisticação dos produtos e um complexo sistema de compensação dos impactos ambientais que causam, isto é, devem ser inovadoras ou intensivas em conhecimento da economia circular \cite{Lucas2019}.

Nesse sentido, a inovação e a tecnologia direcionadas para a sustentabilidade podem ser alternativas para redução de problemas ambientais causados pelo desenvolvimento industrial, que resultou no avanço da produção de novos produtos e serviços, considerando que essa relação entre inovação, tecnologia e sustentabilidade sejam abrangentes, obtendo vários conceitos e definições como inovação sustentável, verde, eco ou ambiental, e tem por objetivo principal a redução de riscos ambientais, poluição e outros impactos negativos ao meio ambiente \cite{Pinsky2017}.

Nos últimos anos, no Brasil, vários trabalhos vêm sendo desenvolvidos no campo das tecnologias sustentáveis com foco na minimização dos impactos ambientais. Estratégias como gerenciamento adequado dos resíduos sólidos e líquidos, uso de fontes energéticas renováveis e produção local dos alimentos vêm sendo cada vez mais desenvolvidas. As técnicas supracitadas auxiliam em um melhor desempenho dos índices sociais, econômicos e ambientais das populações quando relacionadas à sustentabilidade.

Estratégias de marketing e negócios concebidos nos princípios do Desenvolvimento Sustentável e de Tecnologias Apropriadas (TA) poderão ser importantes na definição de tecnologias-chave onde o país deva investir, tanto para a resolução dos seus problemas ambientais básicos quanto para uma política de exportação de produtos, principalmente para os países que presam a sustentabilidade ambiental. 

Para \citeonline{Teixeira2019a}, as políticas públicas — tidas como diretrizes e princípios norteadores de ação do poder público — têm caráter estratégico para a condução de uma nação. Quando associadas à CT\&I, essas políticas são fatores determinantes para o crescimento econômico de um país, especialmente no que diz respeito à capacidade de produção de bens, serviços e a influência no padrão de vida de sua população \cite{Corazza2004}.

As pesquisas em marketing e empreendedorismo devem apresentar soluções visando controlar os problemas socioambientais provocados pelo desenvolvimento predatório e manter o empreendedor economicamente competitivo. Atualmente há um consenso de que a pesquisa é um investimento fundamental para o desenvolvimento sustentável e para a melhoria da qualidade de vida dos povos \cite{Bufrem2018}. 

O desenvolvimento sustentável e a responsabilidade social assumem um papel cada vez mais importante nas estratégias das empresas. O agravamento dos problemas ambientais está ligado à maneira como o conhecimento técnico-científico tem sido aplicado no processo produtivo. Os danos ao meio ambiente não são fatos inesperados e imprevisíveis, apenas demonstram a falta de capacidade do conhecimento de controlar os efeitos gerados pelo desenvolvimento industrial \cite{Maranhao2016}.

Considerando a importância de estudos desta natureza, tanto para a definição de prioridades e fomento de investimento no campo das políticas governamentais relacionadas ao desenvolvimento científico e tecnológico do país quanto para sua implementação e avaliação no nível local, este estudo propõe um modelo empírico-prático para validação de marcas baseadas nos elementos de \textit{Ecobranding} relacionado à sustentabilidade ambiental da Amazônia do Brasil.